\documentclass[11pt,a4paper]{article}
\newcommand{\fichetype}{Exercices}
% Thème A4 commun, librement modifiable.
\usepackage[T1]{fontenc}
\usepackage[utf8]{inputenc}
\usepackage[french]{babel}
\usepackage{lmodern,amsmath,amssymb,mathtools,tikz,pgfplots}
\usepackage[margin=2.2cm,headheight=15pt]{geometry}
\usepackage{xcolor,fancyhdr,enumitem,hyperref,needspace}
\usetikzlibrary{arrows.meta,positioning,calc}
\pgfplotsset{compat=1.18,width=11cm,height=6cm}
\definecolor{Ink}{HTML}{202124}
\definecolor{Blue}{HTML}{356B9E}
\definecolor{Muted}{HTML}{737579}
\color{Ink}
\renewcommand{\familydefault}{\sfdefault}
\setlength{\parindent}{0pt}
\setlength{\parskip}{6pt}
\setlist{itemsep=3pt,topsep=4pt}
\pagestyle{fancy}\fancyhf{}
\fancyhead[L]{\small\color{Muted}MATHS LYCÉE}
\fancyhead[R]{\small\color{Muted}2026--2027}
\fancyfoot[L]{\small\color{Muted}\fichetype}
\fancyfoot[R]{\small\thepage}
\renewcommand{\headrulewidth}{0.3pt}
\hypersetup{colorlinks=true,urlcolor=Blue,linkcolor=Blue}
\newcommand{\R}{\mathbb{R}}
\newcommand{\N}{\mathbb{N}}
\newcommand{\E}{\mathbb{E}}
\newcommand{\V}{\mathrm{V}}
\newcommand{\exercice}[3]{\Needspace{6\baselineskip}\vspace{0.5cm}{\color{Blue}\small\textbf{EXERCICE #1}\hfill #2}\par{\large\bfseries #3}\par}

\begin{document}
{\small\color{Blue}PREMIÈRE SPÉCIALITÉ}\par{\LARGE\bfseries Conditionnement et indépendance}\par{\large Exercices}\par\vspace{0.2cm}\hrule\vspace{0.4cm}
Les exercices progressent de l’application directe vers la recherche. Détailler les calculs et justifier les réponses.
\par\noindent\begin{minipage}{\linewidth}\exercice{1}{Application}{Un tableau à compléter}
Parmi $200$ élèves, $80$ suivent le club sciences $S$, $50$ le club théâtre $T$ et $20$ les deux. Compléter un tableau à double entrée. Calculer $P(S\cap T)$, $P_S(T)$ et $P_T(S)$.
\end{minipage}\par
\par\noindent\begin{minipage}{\linewidth}\exercice{2}{Application}{Lecture d’un arbre}
$P(A)=0,7$, $P_A(B)=0,2$ et $P_{\overline A}(B)=0,6$. Donner les six poids de l’arbre, puis les quatre probabilités des chemins.
\end{minipage}\par
\par\noindent\begin{minipage}{\linewidth}\exercice{3}{Entraînement}{Deux fournisseurs}
Un magasin reçoit $75\%$ de ses produits de $F_1$ et $25\%$ de $F_2$. Les taux de défaut sont $1\%$ et $4\%$. Calculer la probabilité $P(D)$ d’un défaut, puis $P_D(F_2)$.
\end{minipage}\par
\par\noindent\begin{minipage}{\linewidth}\exercice{4}{Entraînement}{Vérifier l’indépendance}
On lance un dé équilibré. $A=\{2,4,6\}$ et $B=\{1,2,3,4\}$. Sont-ils indépendants ? Même question avec $C=\{1,3,5\}$ à la place de $B$.
\end{minipage}\par
\par\noindent\begin{minipage}{\linewidth}\exercice{5}{Entraînement}{Trois tirs}
Un joueur réussit chaque tir avec probabilité $0,7$. On suppose ses trois tirs indépendants. Calculer la probabilité de trois réussites, d’exactement deux réussites et d’au moins une réussite.
\end{minipage}\par
\par\noindent\begin{minipage}{\linewidth}\exercice{6}{Entraînement}{Quatre composants}
Quatre composants fonctionnent indépendamment avec une probabilité $0,9$ chacun. Un dispositif fonctionne s’ils fonctionnent tous. Un autre fonctionne dès qu’au moins un fonctionne. Comparer leurs probabilités de fonctionnement.
\end{minipage}\par
\par\noindent\begin{minipage}{\linewidth}\exercice{7}{Approfondissement}{Trois catégories}
Une enquête fictive distingue trois catégories $A,B,C$ de proportions $0,2$, $0,5$, $0,3$. La probabilité d’une réponse positive $R$ vaut respectivement $0,8$, $0,4$, $0,1$. Calculer $P(R)$, puis la proportion de $A$ parmi les réponses positives. $A$ et $R$ sont-ils indépendants ?
\end{minipage}\par
\par\noindent\begin{minipage}{\linewidth}\exercice{8}{Approfondissement}{Un tirage sans remise}
Une urne contient $3$ boules rouges et $2$ bleues. On tire deux boules sans remise. Construire les probabilités des branches, calculer la probabilité de deux rouges et montrer que les couleurs des deux tirages ne sont pas indépendantes.
\end{minipage}\par
\end{document}
